\chapter{Local centres and boundary generators}
\label{cyp:section}

A representation of a boundary algebra becomes a reconstruction
of the bulk only after its operations determine the bulk action.
The tensor used by the action matters: convolution combines
separated positions, whereas the diagonal tensor places an
interior input above a boundary input at the same position.
Each tensor has its own representing problem. We begin with a
local geometric model in which the boundary data and their
maps can be recovered completely.

On a disk, restriction to the puncture and to the closed point
misses the homotopy attaching them. The exceptional component
of the trace \eqref{cyt:eq:trace} is such a homotopy.
Nontrivial monodromy supplies a second difficulty: every finite
collection of finite-rank probes misses some rank-one local
system. The regular monodromy module solves this problem in
an explicitly enlarged category. Its derived centre then
remembers the attaching point through logarithmic derivatives.

This calculation leads to three further constructions: the
categorical centre of a tensor category, the internal centre
for convolution, and the ambient endomorphism coalgebra for
the diagonal tensor. The last exists without a factorization
condition. Requiring factorizing actors changes the answer:
a nonzero vertex derivation produces a parallel pair that no
receiving factorizing monoid action can identify. The proof
retains the action paths and multiplication coherences.

Let \(D\) be a complex disk, \(U=D\setminus\{0\}\), and
let \(j:U\hookrightarrow D\) and \(i:\{0\}\hookrightarrow D\)
be the inclusions. Fix a field \(k\) of characteristic zero.
Work in the presentable stable category \(\mathcal C(D)\)
of sheaves of \(k\)-complexes whose restriction to \(U\)
is locally constant. Fibers and complexes are unrestricted
in rank and degree.

All tensor products and mapping complexes
below are derived unless a chain model is explicitly given.
Derived categories retain their mapping complexes, not only
their homotopy categories.
Cycle traces are extended from \(\mathbb Q\) to \(k\).

Choose the positively oriented generator \(s\) of
\(\pi_1(U)\), and write \(\Lambda=k[s,s^{-1}]\).
A local system of complexes is determined by a left
\(\Lambda\)-module complex \(M\). Denote it by
\(\mathcal L_M\). The link functor is
\begin{equation}\label{cyp:eq:link}
 K(M)=R\operatorname{Hom}_{\Lambda}(k,M)
        \simeq[M\xrightarrow{s-1}M].
\end{equation}

The link degrees are zero and one, totalized with the
internal differential of \(M\). The free resolution
\(0\to\Lambda\xrightarrow{s-1}\Lambda\to k\to0\)
both proves this formula and supplies a functorial chain
model. In particular \(i^*Rj_*\mathcal L_M=K(M)\).

\section{Objects and maps in the attaching diagram}

\begin{proposition}\label{cyp:prop:gluing}
An object of \(\mathcal C(D)\) is equivalently a diagram
\begin{equation}\label{cyp:eq:data}
           (M,V,\alpha:V\longrightarrow K(M)),
\end{equation}
with \(M\in D(\Lambda)\) and \(V\in D(k)\).
The reconstructed sheaf and its costalk are
\begin{align}
 F&=\operatorname{Fib}\left(
 Rj_*\mathcal L_M\oplus i_*V
 \xrightarrow{\mathrm{res}-i_*\alpha}i_*K(M)\right),
                  \label{cyp:eq:reconstruct}\\
 i^!F&=\operatorname{Fib}\bigl(V\xrightarrow{\alpha}K(M)\bigr).
                  \label{cyp:eq:costalk}
\end{align}
For two such diagrams, the mapping complex is
\begin{equation}\label{cyp:eq:maps}
 \operatorname{Fib}\left(
 R\operatorname{Hom}_{\Lambda}(M,M')\oplus
 R\operatorname{Hom}_k(V,V')
 \xrightarrow{\Phi}R\operatorname{Hom}_k(V,K(M'))\right),
\end{equation}
where \(\Phi(a,b)=K(a)\alpha-\alpha'b\).
\end{proposition}

\begin{proof}
Restricting \eqref{cyp:eq:reconstruct} to \(U\) gives
\(\mathcal L_M\). At the point, the fiber of
\(K(M)\oplus V\xrightarrow{(1,-\alpha)}K(M)\)
is \(V\), with the prescribed attaching map \(\alpha\).
Conversely, the adjunction maps from any \(F\) to
\(Rj_*j^*F\) and to \(i_*i^*F\) define a map to this
fiber. Its restrictions to both strata are equivalences.
The open--closed localization triangle shows it is an
equivalence. Applying that triangle at the closed stratum
gives \eqref{cyp:eq:costalk}.

A map of the reconstructed diagrams consists of maps on
the two vertices and a homotopy commuting the attaching
square. The homotopy limit of these mapping complexes is
exactly \eqref{cyp:eq:maps}. This computes maps in the
derived diagram category, including their higher homotopies.
The reconstruction equivalence transfers it to
\(\mathcal C(D)\).
\end{proof}

Choose compatible resolutions so that \(\alpha,\alpha'\)
and \(\Phi\) are actual chain maps. In the fiber convention
used here, a degree-\(n\) element is \((a,b,h)\), where
\(a,b\) have degree \(n\) and \(h\) degree \(n-1\), with
\begin{equation}\label{cyp:eq:differential}
 d(a,b,h)=(da,db,K(a)\alpha-\alpha'b-dh).
\end{equation}
Thus a degree-zero map with fixed closed endpoints exists
precisely when \(K(a)\alpha-\alpha'b\) is a boundary.
Its homotopy component is an actual primitive in that
mapping complex. The supported extension problem below uses the same
homotopy-fiber construction.

\begin{proposition}\label{cyp:prop:gysin}
In \eqref{cyp:eq:maps}, the oriented Gysin morphism
\(\gamma:i_*k[-2]\to k_D\) has representative
\((0,0,h_\gamma)\), where \(h_\gamma\) sends the degree-two
point generator to the degree-one link generator.
This class spans the degree-zero cohomology of the mapping
complex. With these conventions the normalized cycle trace
of Theorem~\ref{cyt:thm:cycle-trace} has open component
\(a=\operatorname{id}\), closed component
\(b=(\operatorname{id},0):k\oplus k[-2]\to k\),
and homotopy component \(\tfrac13h_\gamma\) on the second
summand.
\end{proposition}

\begin{proof}
The constant sheaf has diagram \((k,k,\alpha)\), where
\(K(k)=[k\xrightarrow0 k]\) and \(\alpha\) is inclusion
into link degree zero. The point source has diagram
\((0,k[-2],0)\).

Its mapping complex has a \(b\)-component in degree \(-2\).
The \(h\)-components from link degrees zero and one have
degrees \(-1\) and zero, respectively. Formula~\eqref{cyp:eq:differential} makes the
first differential \(-\operatorname{id}\). The second is
zero. Exactly the asserted class survives.

The costalk fiber differential is \(d(v,u)=(dv,\alpha v-du)\).
A closed triple \((a,b,h)\) induces the costalk map
\((v,u)\mapsto(bv,K(a)u-hv)\).

Let \(\eta\) be the positive link class with integral one.
On a small closed subdisk \(D_\epsilon\), relative
integration is \(\int_{D_\epsilon}v-\int_{\partial D_\epsilon}u\).
Thus \((0,-\eta)\) is the positive degree-two orientation
class in the relative cone.

Consequently \(h_\gamma(1)=\eta\) induces the identity on the
oriented costalk \(k[-2]\). This is precisely the Gysin
counit, also normalized by \(H_c^2(\gamma)=1\).
Adding the identity on the constant summand and multiplying
this class by \(1/3\) gives the last assertion.
\end{proof}

This calculation locates the lost datum. The exceptional
ordinary stalk map is zero, but its attaching homotopy is
nonzero. Extraordinary costalks and compact support already
detect its coefficient, as in Theorem~\ref{cyf:thm:gysin}.

\section{Probes and monodromy}

\begin{definition}
For \(P,F\in\mathcal C(D)\), the \emph{probe value} of
\(F\) at \(P\) is \(R\operatorname{Hom}(P,F)\).
A family of probes is \emph{conservative} if their values
vanish simultaneously only on the zero object.
\end{definition}

The preceding Gysin calculation uses constant monodromy.
Nontrivial monodromy can make both constant-probe values
vanish, even for a rank-one local system.

For \(\lambda\in k^\times\), let \(k_\lambda\) be the
one-dimensional \(\Lambda\)-module on which \(s\) acts
by \(\lambda\), and write \(\mathcal L_\lambda\) for its
local system.

\begin{proposition}\label{cyp:prop:blindness}
If \(\lambda\ne1\), the nonzero object
\(F_\lambda=j_!\mathcal L_\lambda\) satisfies
\[
 F_\lambda\simeq Rj_*\mathcal L_\lambda,\qquad
 R\operatorname{Hom}(k_D,F_\lambda)=0,
 \qquad R\operatorname{Hom}(i_*k,F_\lambda)=0,
 \qquad R\Gamma_c(D,F_\lambda)=0.
\]
More generally, a finite family of bounded finite-rank
constructible objects cannot be conservative as probes
\(R\operatorname{Hom}(P,-)\) on all bounded finite-rank
objects of \(\mathcal C(D)\).
\end{proposition}

\begin{proof}
The link complex of \(k_\lambda\) has differential
\(\lambda-1\), hence is acyclic. The localization triangle
identifies extension by zero with derived direct image,
and \eqref{cyp:eq:costalk} makes its costalk zero.
Ordinary sections are the same acyclic circle complex.
Compact sections of \(U\simeq S^1\times(0,1)\) are
that complex tensored with the compact cochains of the
open interval, so vanish as well.

For finitely many bounded probes, choose \(\lambda\ne0,1\)
outside the spectra of all their finite-dimensional open
cohomology modules. There are only finitely many excluded
values, and \(k\) is infinite. If \(W\) is a finite-dimensional
monodromy module with operator \(T_W\), it has free resolution
\[
 0\longrightarrow\Lambda\otimes_k W
 \xrightarrow{s\otimes1-1\otimes T_W}\Lambda\otimes_k W
 \longrightarrow W\longrightarrow0.
\]

Injectivity follows by comparing the highest Laurent power.
The cokernel imposes exactly the given monodromy relation.
Applying \(\operatorname{Hom}_{\Lambda}(-,k_\lambda)\)
gives an invertible differential when \(\lambda\) is outside
the spectrum. Truncation triangles give the same vanishing
for each bounded open complex.

In \eqref{cyp:eq:maps},
the target has zero closed term and acyclic link, so these
open vanishings prove \(R\operatorname{Hom}(P,F_\lambda)=0\)
for every chosen probe. Its restriction to \(U\) is a
nonzero rank-one local system.
\end{proof}

\subsection{The ramified double cover}

The monodromy \(-1\) occurs in a branched correspondence.
Let \(\pi:\widetilde D\to D\) be the squaring map
\(z=u^2\) between unit disks, and let \(\sigma(u)=-u\).

\begin{example}\label{cyp:ex:ramification}
The coefficient object and its complex-oriented trace are
\[
 R\pi_*k_{\widetilde D}=k_D\oplus j_!\mathcal L_{-1},
 \qquad \operatorname{Tr}_\pi=(2,0),
 \qquad j_!\mathcal L_{-1}=Rj_*\mathcal L_{-1}.
\]

Indeed, the idempotents \((1\pm\sigma)/2\) split the
permutation representation on the two sheets. At zero
there is one invariant stalk. The anti-invariant link has
differential \(-2\), so its stalk and costalk vanish.
The finite map has no higher direct images. On the
unramified open the oriented trace sums the two sheets.
It extends as multiplication by two on the constant
summand, including the ramification point.

This is the normal form of the first nested Hilbert
correspondence on \(\mathbb C^2\). After centering a length-two
quotient, a chart with basis \(1,X\) has equations
\[
                       Y-rX=0,\qquad X^2=z/4.
\]
Here \(r\) fixes its tangent direction. Choosing a retained
length-one quotient means choosing the root \(X=-u/2\),
so the forgetful map is \(z=u^2\). The center and \(r\)
are smooth parameters. Thus the normal anti-invariant
coefficient is exactly \(\mathcal L_{-1}\).
\end{example}

This coefficient is missed by the two constant probes.
A generator for all monodromies must contain more information.

\section{The regular monodromy generator}

The regular module \(\Lambda\) supplies all monodromies at
once. Its local system has infinite rank. No nonzero Laurent
polynomial lies in a finite-dimensional \(s\)-stable subspace:
its translates by increasing powers of \(s\) are linearly
independent. Thus \(\mathcal L_\Lambda\) is not a filtered
union of finite-rank local systems. The unrestricted category
specified above is essential.

Set
\begin{equation}\label{cyp:eq:generator}
 A=j_!\mathcal L_\Lambda,\qquad B=i_*k,
 \qquad \mathcal G=A\oplus B.
\end{equation}
For a sheaf with diagram \((M,V,\alpha)\), adjunction gives
\begin{equation}\label{cyp:eq:probe-values}
 R\operatorname{Hom}(A,F)=M,
 \qquad R\operatorname{Hom}(B,F)=i^!F.
\end{equation}

\begin{definition}
An object \(P\) of \(\mathcal C(D)\) is \emph{compact}
if \(R\operatorname{Hom}(P,-)\) commutes with arbitrary
sums. A family \emph{generates} if the smallest stable
subcategory containing it and closed under arbitrary sums
is \(\mathcal C(D)\).
\end{definition}

\begin{theorem}\label{cyp:thm:generator}
The object \(\mathcal G\) is a compact generator of
\(\mathcal C(D)\). Its derived endomorphism algebra has the
triangular model
\begin{equation}\label{cyp:eq:endomorphisms}
 \mathcal E=
 \begin{pmatrix}\Lambda&k[-2]\\0&k\end{pmatrix},
 \qquad d=0,
\end{equation}
where the left action of \(\Lambda\) on the off-diagonal
entry is evaluation at \(s=1\). Consequently
\(\mathcal C(D)\) is equivalent to the derived category
of right \(\mathcal E\)-modules, retaining their full
derived mapping complexes.
\end{theorem}

\begin{proof}
The first functor in \eqref{cyp:eq:probe-values} commutes
with arbitrary sums. The second does too: by
\eqref{cyp:eq:costalk} it is a fiber of stalk and link
functors, and \eqref{cyp:eq:link} is a finite free
resolution. Hence \(A,B\) are compact.

If both probe
values vanish, the open restriction is zero, so localization
makes \(F=i_*V\). Its costalk is then \(V\), which is zero.
More explicitly, \(\Lambda\) generates \(D(\Lambda)\)
and \(k\) generates \(D(k)\). The triangle
\(j_!j^*F\to F\to i_*i^*F\) puts every \(F\) in the
localizing subcategory generated by \(A,B\).

Use their diagrams \((\Lambda,0,0)\) and \((0,k,0)\)
in \eqref{cyp:eq:maps}. The four blocks are
\begin{align*}
 R\operatorname{Hom}(A,A)&=\Lambda,&
 R\operatorname{Hom}(A,B)&=0,\\
 R\operatorname{Hom}(B,B)&=k,&
 R\operatorname{Hom}(B,A)&=K(\Lambda)[-1].
\end{align*}
Let \(Q\) be the complex with \(Q^1=Q^2=\Lambda\)
and differential \(s-1\). It identifies with
\(K(\Lambda)[-1]\) by multiplying degree one by \(-1\).
Composition of diagrams therefore gives a strict model
\[
 \widetilde{\mathcal E}
       =\begin{pmatrix}\Lambda&Q\\0&k\end{pmatrix}.
\]

Evaluation at one in degree two and zero in degree one
give a quasi-isomorphism of bimodules \(Q\to k[-2]\).
It extends to a dg algebra quasi-isomorphism
\(\widetilde{\mathcal E}\to\mathcal E\): two off-diagonal
entries cannot compose, and all remaining products are
the indicated left and right module actions.

The functor \(R\operatorname{Hom}(\mathcal G,-)\) takes
values in right \(\widetilde{\mathcal E}\)-modules by
precomposition. Its left adjoint is
\(N\mapsto N\otimes^L_{\widetilde{\mathcal E}}\mathcal G\).
The unit and counit are equivalences on the two idempotent
free modules and the corresponding summands of
\(\mathcal G\). Both functors preserve sums and triangles,
so the equivalences extend to their generated categories,
which are the entire categories just specified.

Finally transport modules along
\(\widetilde{\mathcal E}\to\mathcal E\). If an explicit
inverse over the zero-differential algebra is desired, put
\(\mathcal G_{\mathcal E}
=\mathcal E\otimes^L_{\widetilde{\mathcal E}}\mathcal G\).
The inverse is
\(N\mapsto N\otimes^L_{\mathcal E}\mathcal G_{\mathcal E}\).
Derived transport defines the left action without factoring
the original chain action strictly through the quotient.
\end{proof}

\begin{example}\label{cyp:ex:extension}
Full derived module maps carry more than their underlying
graded cohomology maps. Fix \(\lambda\ne0,1\). The sequence
\[
 0\longrightarrow k_\lambda\longrightarrow
 \Lambda/(s-\lambda)^2\longrightarrow k_\lambda
 \longrightarrow0
\]
does not split as \(\Lambda\)-modules. Indeed, any lift of
the quotient generator is sent by \(s-\lambda\) to the
nonzero submodule generator. Its connecting morphism is a
nonzero element of
\(\operatorname{Ext}^1_\Lambda(k_\lambda,k_\lambda)\).

Since \(j_!\) is fully faithful, it gives a nonzero map
\(F_\lambda\to F_\lambda[1]\).
The underlying map of \(A\)-probe complexes is zero in
\(D(k)\), because the sequence splits as vector spaces.
Both \(B\)-probe complexes are zero. The derived module
map in Theorem~\ref{cyp:thm:generator} nevertheless remains
nonzero. Conservativity on objects cannot be replaced by
faithfulness of the underlying graded vector-space maps.
\end{example}

\section{Independent normal directions}

Give \(D^r\) the coordinate stratification, and allow arbitrary
local systems of complexes on each stratum. For each
\(I\subseteq\{1,\ldots,r\}\), take the external product
which uses \(A\) in coordinates in \(I\) and \(B\) in
the others.

\begin{proposition}\label{cyp:prop:normal-crossings}
These \(2^r\) external products form a compact generating
family. The derived endomorphism algebra of their sum is
quasi-isomorphic to the graded tensor product
\(\mathcal E^{\otimes r}\), with its Koszul signs.
It reconstructs all objects and maps in this stratified
category. In particular it detects a nonzero discrepancy
from the two-axis coefficient \(1/9\) of
Corollary~\ref{cyt:cor:products}.
\end{proposition}

\begin{proof}
Apply open--closed localization one coordinate at a time.
On a stratum with open coordinates \(I\), the regular
local system is the free module over
\(k[s_i^{\pm1}:i\in I]\). It generates that stratum's
local systems. The localization triangles give generation
on the whole polydisk. Compactness follows from the finite
tensor product of the two-term link resolutions.

The same resolutions compute each Hom block between the
displayed generators. A factor with source \(A\) and
target \(B\) gives zero. The other factors give
\(\Lambda\), \(k\), or \(Q\), as in the preceding proof.
The product cellular complex is the tensor product of
these finite free resolutions. Composition is tensor
composition with the cohomological signs. Thus its full
dg endomorphism algebra is
\(\widetilde{\mathcal E}^{\otimes r}\).

Tensoring the block quasi-isomorphisms over the field \(k\)
gives \(\mathcal E^{\otimes r}\), and the same adjunction
argument reconstructs the category. The exceptional
two-axis Gysin class is nonzero by coordinate purity,
so a nonzero scalar multiple remains a nonzero derived
module map under this equivalence.
\end{proof}

For independent flags this model tests the whole normalized
trace, including its attaching homotopies. An interacting
flag requires its own stratification and link functors.
The reconstruction theorem applies once those coefficient
objects and their attaching maps have been supplied.

\section{Derived centres and change of generator}
\label{cycen:sec:derived}

A compact generator reconstructs objects and maps from one
algebra. The operations that act naturally on every module
are its derived centre. This requires the bimodule structure
of the algebra itself.

For a unital differential graded algebra $\mathsf A$ over $k$,
put $\mathsf A^e=\mathsf A\otimes_k\mathsf A^{\mathrm{op}}$ and define
\[
 Z_{\mathrm{dg}}(\mathsf A)
   =R\operatorname{Hom}_{\mathsf A^e}(\mathsf A,\mathsf A),
 \qquad
 HH^n(\mathsf A)=H^n Z_{\mathrm{dg}}(\mathsf A).
\]
The ordinary centre is $HH^0$ when $\mathsf A$ is concentrated
in degree zero. The whole complex has a product and coherent
operations refining its commutativity.

Write $s=[1]$ for suspension, so $|sa|=|a|-1$.
The normalized two-sided bar resolution gives the product
complex
\begin{equation}\label{cycen:eq:bar}
 C(\mathsf A)=\prod_{q\geq0}
  \operatorname{Hom}_k\bigl((s\overline{\mathsf A})^{\otimes q},
                                  \mathsf A\bigr),
 \qquad \overline{\mathsf A}=\mathsf A/k\mathbf1.
\end{equation}
Derived representatives are understood. An arity-$q$ cochain
vanishes if any of its inputs is the unit. The product over
arities is part of the totalization; it is not replaced by a
direct sum when infinitely many arities contribute in one degree.

The signs can be specified on suspended inputs. For a
homogeneous cochain $f$, set $F=sf$, with $|F|=|f|-1$.
The brace $F\{G_1,\ldots,G_r\}$ sums over ordered disjoint
insertions of the $G_j$ into $F$. On suspended inputs $x_i$,
an insertion beginning at positions $\ell_j$ has sign
\begin{equation}\label{cycen:eq:brace-sign}
 (-1)^{\sum_j|G_j|\sum_{i<\ell_j}|x_i|}.
\end{equation}
An arity-zero insertion consumes no input. The inner sum
includes inputs consumed by earlier inserted maps. Projection
to $s\overline{\mathsf A}$ precedes every insertion into a
normalized input.

Let $b=b_1+b_2$ be the bar coderivation on $T^c(s\mathsf A)$,
with
\[
 b_1(sa)=-s(da),\qquad b_2(sa,sb)=(-1)^{|a|}s(ab).
\]
For suspended maps put
$[F,G]=F\{G\}-(-1)^{|F||G|}G\{F\}$.
The cochain differential, cup product and Gerstenhaber bracket
are
\begin{equation}\label{cycen:eq:operations}
 Df=-s^{-1}[b,F],\qquad
 f\smile g=(-1)^{|f|}s^{-1}b_2\{F,G\},\qquad
 [f,g]_{\mathrm G}=s^{-1}[F,G].
\end{equation}
The bracket has degree $-1$.
An ordinary arity-$q$ multilinear map is put in this
convention by precomposing with $(s^{-1})^{\otimes q}$,
including its tensor Koszul sign.

\begin{proposition}\label{cycen:prop:bar-operations}
The complex \eqref{cycen:eq:bar} computes
$Z_{\mathrm{dg}}(\mathsf A)$. Its cup product represents
composition of derived bimodule endomorphisms. Its cup and
braces give the Hochschild $E_2$ structure. On cohomology,
\eqref{cycen:eq:operations} defines a Gerstenhaber algebra.
\end{proposition}
\begin{proof}
The two-sided bar is augmented to $\mathsf A$ by multiplication.
Inserting the unit contracts the augmented resolution. Applying
bimodule Hom produces \eqref{cycen:eq:bar}. Its differential
has the internal differential, the adjacent products and the
two outer actions, exactly the terms of $-s^{-1}[b,F]$.

The Alexander--Whitney diagonal cuts a bar string at every
position. Applying $g$ to its final segment and then $f$ to
the preceding segment gives $f\smile g$, with the sign of
moving the maps through the suspended inputs. The diagonal
commutes with the bar differential: products across a cut
cancel against its outer actions. Its coassociativity identifies
cup with composition of the resulting chain endomorphisms.

A second family of brace insertions either lies inside an
existing insertion or occupies a remaining slot. These choices
partition the same ordered substitutions, with the same Koszul
signs. This proves the brace identities; $b^2=0$ gives their
differential identities. The operad generated by cup, unit and
these braces is a chain model of the little two-cubes operad
\cite{McClureSmith}*{Theorem~1.1}. Its cohomology operations
are the asserted product and bracket.
\end{proof}

\begin{example}\label{cycen:ex:dual-numbers}
Let $\mathsf A=k[\epsilon]/(\epsilon^2)$. The periodic
bimodule resolution has successive differentials alternating
between multiplication by $x-y$ and $x+y$ in
$\mathsf A^e=k[x,y]/(x^2,y^2)$, beginning with $x-y$.
The kernels are the images of the other multiplication:
this follows by evaluating on the basis $1,x,y,xy$.
After applying bimodule Hom, the differentials alternate
between zero and multiplication by $2\epsilon$. Hence
\[
 HH^0(\mathsf A)=\mathsf A,\qquad
 HH^{2r+1}(\mathsf A)=k\epsilon\quad(r\geq0),\qquad
 HH^{2r}(\mathsf A)=k\quad(r\geq1).
\]
In contrast, $\operatorname{End}_k(\mathsf A)$ is
four-dimensional and concentrated in degree zero. Linear
endomorphisms do not compute the derived centre.
\end{example}

\begin{theorem}\label{cycen:thm:morita}
A $k$-linear equivalence of presentable derived module
categories, preserving colimits, identifies their derived
centres as $E_2$ algebras and intertwines their actions on
modules. In particular, changing a compact generator preserves
the full Hochschild structure.
\end{theorem}
\begin{proof}
A colimit-preserving $k$-linear functor from left
$\mathsf A$-modules to left $\mathsf B$-modules is tensoring
with a $\mathsf B$--$\mathsf A$ bimodule $P$. Indeed, take
$P$ to be its value on $\mathsf A$; the right action comes
from the right action on this free module. The assertion
holds on free modules by linearity and on every module by
its free bar resolution. The same construction identifies
natural transformations and their higher homotopies.

Thus endofunctors of the module category correspond to
bimodules, composition corresponds to balanced tensor, and
the identity corresponds to $\mathsf A$. Its endomorphism
complex is $Z_{\mathrm{dg}}(\mathsf A)$. Vertical composition
cuts bar strings; horizontal composition substitutes their
components. These are cup and braces. Conjugating
endofunctors by the given equivalence and a chosen inverse
preserves both compositions and their interchange. It therefore
identifies the $E_2$ algebras and their evaluations. This is
also the derived invariance of Hochschild higher operations
\cite{cycen:KellerHigher}*{Sections~3.2 and~4.6}.
\end{proof}

For $\mathsf B=M_n(\mathsf A)$ the comparison is visible
entry by entry. With $e=E_{11}$, the bimodules
$P=\mathsf Be$ and $Q=e\mathsf B$ satisfy
\[
 Q\otimes_{\mathsf B}^{\mathbb L}P\simeq\mathsf A,
 \qquad P\otimes_{\mathsf A}^{\mathbb L}Q\simeq\mathsf B.
\]
For an unsuspended arity-$q$ cochain,
\begin{equation}\label{cycen:eq:matrix}
 (\Phi_nf)(X_1,\ldots,X_q)_{i_0i_q}
  =\sum_{i_1,\ldots,i_{q-1}}
   f\bigl((X_1)_{i_0i_1},\ldots,(X_q)_{i_{q-1}i_q}\bigr).
\end{equation}
At arity zero it is the diagonal matrix with entry $f$.
Concatenating index paths proves compatibility with cup;
substituting paths proves compatibility with every brace.
The relative bar over $M_n(k)$ computes the ordinary bar
Hom because this algebra is separable: inserting
$\sum_i E_{i1}\otimes E_{1i}$ contracts the extra balancing
directions. Its relative cochains identify with $C(\mathsf A)$
by \eqref{cycen:eq:matrix}. Thus the displayed map is a
quasi-isomorphism with all these operations.

\section{The centre retains the attaching point}
\label{cycen:sec:attached}

Return to the disk generator $A\oplus B$ of
\eqref{cyp:eq:generator}. Replace it by $A\oplus B[-2]$.
The off-diagonal derived Hom is now in degree zero, so its
endomorphism algebra is quasi-isomorphic to
\begin{equation}\label{cycen:eq:simple}
 T_x=\begin{pmatrix}\Lambda&k\\0&k\end{pmatrix},
 \qquad \Lambda=k[x,(1+x)^{-1}],\quad x=s-1.
\end{equation}
The arrow $a$ in the upper-right entry satisfies $xa=0$.
The shift changes the compact generator; it does not discard
the degree-two attachment. More explicitly, a dg algebra with
cohomology only in degree zero has the multiplicative truncation
zigzag $E\leftarrow\tau_{\leq0}E\to H^0(E)$.
Theorem~\ref{cycen:thm:morita} therefore identifies the derived
centre of the disk category with that of $T_x$.

To vary the attachment, let $F\in k[x]$ be monic and
nonconstant, with $F(-1)\ne0$, and put
\begin{equation}\label{cycen:eq:attached-family}
 M_F=\Lambda/(F),\qquad
 T_F=\begin{pmatrix}\Lambda&M_F\\0&k\end{pmatrix}.
\end{equation}
Its diagonal idempotents are $e_1,e_2$; the class of $1$ in
$M_F$ is $a=e_1ae_2$, and $F(x)a=0$. The algebra $T_{x^r}$
for $r>1$ defines a different attaching category. A geometric
realization of that category requires additional data.

\begin{theorem}\label{cycen:thm:resolution}
The algebra $T=T_F$ has the projective bimodule resolution
\[
 0\longrightarrow P_2\xrightarrow{d_2}P_1
 \xrightarrow{d_1}P_0\xrightarrow{\mu}T\longrightarrow0,
\]
where
\[
 \begin{aligned}
 P_0&=(Te_1\otimes e_1T)[e_1]\oplus(Te_2\otimes e_2T)[e_2],\\
 P_1&=(Te_1\otimes e_1T)[x]\oplus(Te_1\otimes e_2T)[a],\\
 P_2&=(Te_1\otimes e_2T)[R].
 \end{aligned}
\]
Here all tensors are over $k$, and brackets name generators.
Writing $F(x)=\sum_j f_jx^j$, the differentials are
\begin{equation}\label{cycen:eq:resolution-differentials}
 \begin{aligned}
 d_1[x]&=x[e_1]-[e_1]x,&
 d_1[a]&=a[e_2]-[e_1]a,\\
 d_2[R]&=\sum_{j\geq1}f_j\sum_{b=0}^{j-1}
             x^b[x]x^{j-1-b}a+F(x)[a].
 \end{aligned}
\end{equation}
\end{theorem}
\begin{proof}
Every summand is an idempotent summand of a free $T^e$-module.
The composite $d_1d_2$ telescopes to
$F(x)a[e_2]-[e_1]F(x)a=0$. Check exactness in the four
idempotent blocks. The $(2,1)$ block is zero, the $(2,2)$
augmentation is the identity, and the $(1,1)$ block is the
diagonal resolution of $\Lambda$, whose kernel is generated
by the nonzero divisor $x-y$ in $\Lambda\otimes_k\Lambda$.

For the $(1,2)$ block set
\[
 L=\Lambda_x\otimes_k M_{F,y},\qquad
 D_F(x,y)=\frac{F(x)-F(y)}{x-y}.
\]
The augmented complex is
\[
 0\longrightarrow\Lambda_x\longrightarrow L\oplus\Lambda_x
 \longrightarrow L\oplus M_F\longrightarrow M_F\longrightarrow0
\]
with maps
\[
 W\longmapsto(D_FW,FW),\quad
 (U,V)\longmapsto((x-y)U-V,\overline V),\quad
 (Z,m)\longmapsto Z(y,y)+m.
\]
The first is injective because $F$ is nonzero in the domain
$\Lambda_x$. Multiplication by $x-y$ on $L$ is injective:
it is monic in $x$ before localization, even when $M_F$ is
nonreduced. If $(x-y)U=V$ and $\overline V=0$, write
$V=FW$; then $U=D_FW$. If $Z(y,y)+m=0$, choose a lift
$V\in\Lambda$ of $m$. The element $Z+V$ is divisible by
$x-y$, and its quotient gives a preimage of $(Z,m)$.
The final map is surjective. These arguments prove exactness.
\end{proof}

Let $S=ke_1\oplus ke_2$. The relative normalized bar cochains
\[
 C_S^n(T_F,T_F)=\operatorname{Hom}_{S\text{-}S}
       ((T_F/S)^{\otimes_S n},T_F)
\]
retain cup and braces and compute the derived centre, since
$S$ is separable. The finite resolution maps to this bar by
sending $[x]$ and $[a]$ to their one-bar classes and
\begin{equation}\label{cycen:eq:relation-bar}
 [R]\longmapsto
 \sum_{j\geq1}f_j\sum_{b=0}^{j-1}x^b[x\mid x^{j-1-b}a].
\end{equation}
Its bar boundary telescopes to
\eqref{cycen:eq:resolution-differentials}; the relation
$F(x)a=0$ supplies the constant term. The short complex
below computes cohomology. Its length does not define strict
brace operations on it.

\begin{theorem}\label{cycen:thm:attached-centre}
The underlying derived-centre complex is computed by
\begin{equation}\label{cycen:eq:small}
 \Lambda\oplus k\xrightarrow{d^0}\Lambda\oplus M_F
 \xrightarrow{d^1}M_F,
 \quad
 d^0(f,c)=(0,c-\overline f),\qquad
 d^1(v,w)=\overline{F'v}.
\end{equation}
These terms are in degrees zero, one and two. If
$g=\gcd(F,F')$ is monic, then
\begin{equation}\label{cycen:eq:attached-hh}
 \begin{aligned}
 HH^0(T_F)&=k+F\Lambda,&
 HH^1(T_F)&=(F/g)\Lambda\,\partial_x,\\
 HH^2(T_F)&=\Lambda/(F,F'),& HH^{n}(T_F)&=0\quad(n\geq3).
 \end{aligned}
\end{equation}
The function $f=c+Fh$ represents the central matrix
$\operatorname{diag}(f,c)$.
\end{theorem}
\begin{proof}
Bimodule Hom sends the vertex generators to $(f,c)$ and
the loop and arrow generators to $(v,w)$. Evaluating their
coboundaries on \eqref{cycen:eq:resolution-differentials}
gives $d^0$ and $F'v+Fw=F'v$ modulo $F$. The first
differential surjects onto the arrow component. The remaining
kernel is $\{v:F\mid F'v\}=(F/g)\Lambda$, and the cokernel
is $\Lambda/(F,F')$. The resolution excludes higher groups.
Represent a degree-one class by the derivation
$\delta_v(x)=v$, $\delta_v(a)=\delta_v(e_i)=0$.
The divisibility condition is exactly the condition preserving
$F(x)a=0$.
\end{proof}

For the simple attachment this gives
\begin{equation}\label{cycen:eq:logarithmic-centre}
 HH^0(T_x)=\Lambda,\qquad
 HH^1(T_x)=x\Lambda\partial_x,\qquad HH^{\geq2}(T_x)=0.
\end{equation}
The contraction of \eqref{cycen:eq:small} onto these groups
is explicit:
\[
 \begin{array}{ll}
 i^0(f)=(f,f(0)),&p^0(f,c)=f,\\
 i^1(v)=(v,0),&p^1(v,w)=v-v(0),
 \end{array}
 \qquad
 h^1(v,w)=(0,w),\quad h^2(z)=(z,0).
\]
It satisfies $dh+hd=1-ip$ and $pi=1$. The cup action is
multiplication of functions on vector fields. The derivation
representatives give
\begin{equation}\label{cycen:eq:logarithmic-bracket}
 [v\partial_x,f]=vf',\qquad
 [v\partial_x,w\partial_x]=(vw'-wv')\partial_x,
 \qquad v,w\in x\Lambda.
\end{equation}
Products of two positive-degree classes vanish by
\eqref{cycen:eq:logarithmic-centre}.

\begin{corollary}\label{cycen:cor:missing-derivative}
There is no Gerstenhaber-algebra map
$HH^\bullet(\Lambda)\to HH^\bullet(T_x)$ that is the identity
on degree-zero functions. Consequently there is no unital
$E_2$ map of derived centres with that property.
\end{corollary}
\begin{proof}
The image of $\partial_x$ would be $v\partial_x$ for
$v\in x\Lambda$. Preservation of $[\partial_x,x]=1$
would force $v=1$, whereas $x$ is not invertible in
$\Lambda$. No section property is needed for this contradiction.
\end{proof}

The coordinate $x=s-1$ measures monodromy relative to the
attaching point. It is not a spacetime insertion coordinate.
Restriction to the open object does give a higher Hochschild
map in the opposite direction, by restriction of bar cochains
to a full subcategory \cite{cycen:KellerHigher}*{Section~4.3}.
On cohomology it is the identity on functions and the
inclusion $x\Lambda\partial_x\subset\Lambda\partial_x$.

\subsection{Several attachments and flat deformations}

For $d$ independent simple attachments, tensor the finite
bimodule resolutions. Their terms are finitely generated
projectives, so Hom commutes with the tensors in this
calculation. K\"unneth gives
\begin{equation}\label{cycen:eq:log-polyvectors}
 HH^\bullet(T_x^{\otimes d})=
 k[x_i,(1+x_i)^{-1}:1\leq i\leq d]
   \otimes\Lambda(\theta_1,\ldots,\theta_d),
 \qquad \theta_i=x_i\partial_{x_i},\quad |\theta_i|=1.
\end{equation}
The brackets obey $[\theta_i,x_j]=\delta_{ij}x_i$ and
$[\theta_i,\theta_j]=0$, extended by the graded Leibniz rule.
In two directions, $\theta_1\smile\theta_2\ne0$.

Let $D_i$ be the derivation of $T_x^{\otimes2}$ which acts
by $x_i\partial_{x_i}$ on its open coordinate and kills
both idempotents and both arrows. It preserves the arrow
relations and the localized functions. The commuting
operators $D_1,D_2$ define
\begin{equation}\label{cycen:eq:log-twist}
 a\star_\hbar b=\sum_{n\geq0}\frac{\hbar^n}{n!}
                        D_1^n(a)D_2^n(b).
\end{equation}
This is an associative unital product on
$T_x^{\otimes2}[[\hbar]]$. Both triple products equal ordinary
triple multiplication after
\[
 \exp\bigl(\hbar(D_1\otimes D_2\otimes1
                  +D_1\otimes1\otimes D_2
                  +1\otimes D_1\otimes D_2)\bigr).
\]
The three operators commute, proving associativity in every
formal degree. The derivations kill the unit. The first-order
class is $\theta_1\smile\theta_2$, and
\[
 x_1\star_\hbar x_2=e^\hbar x_1x_2,
 \qquad x_2\star_\hbar x_1=x_1x_2.
\]
The antisymmetrization of a Hochschild coboundary on the
central pair $(x_1,x_2)$ is zero, whereas that of the displayed
first-order term is $x_1x_2$. This also detects its nonzero class.

For a thick attachment $F=x^r$, $r\geq2$, the groups are
\[
 HH^0(T_{x^r})=k+x^r\Lambda,\quad
 HH^1(T_{x^r})=x\Lambda\partial_x,\quad
 HH^2(T_{x^r})=k[x]/(x^{r-1}).
\]
The degree-two directions distinguish it from the simple
attachment even up to derived Morita equivalence. For
$v=xw$ their Gerstenhaber action is
\begin{equation}\label{cycen:eq:thick-bracket}
 [xw\partial_x,g]=xwg'-rwg\pmod{x^{r-1}}.
\end{equation}
Indeed, $x\mapsto x+\varepsilon xw$ carries $x^r$ to
$(1+\varepsilon rw)x^r$. Apply this change to $x^r-\tau g$
and divide the relation by that unit. Its new deformation
term is $g+\varepsilon(xwg'-rwg)$, proving the formula.
In particular $[x\partial_x,x^j]=(j-r)x^j$ in degree two.

All these directions integrate together. Put
\[
 R=k[[t_0,\ldots,t_{r-2}]],\qquad
 F_{\mathbf t}(x)=x^r-\sum_{j=0}^{r-2}t_jx^j,
\]
and complete in the maximal ideal of $R$. Then
\begin{equation}\label{cycen:eq:formal-attaching}
 \mathcal T_{\mathbf t}=
 \begin{pmatrix}
 \Lambda\widehat\otimes_k R&R[x]/(F_{\mathbf t})\\0&R
 \end{pmatrix}
\end{equation}
is a formally flat algebra. The arrow module is free over
$R$ on $1,x,\ldots,x^{r-1}$, and $1+x$ is invertible on
it because $F_{\mathbf t}(-1)$ is a unit. Each Artinian
reduction is free over its parameter ring. The $t_j$ tangent
class is $+x^j$: in the companion presentation, the first
product variation sends $(x,x^{r-1}a)$ to $x^ja$, and
\eqref{cycen:eq:relation-bar} identifies this with the
positive relation class.

For $r=2$, multiplication by $1+x$ on the arrow module is
\[
 S_\tau=\begin{pmatrix}1&\tau\\1&1\end{pmatrix},\qquad
 (S_\tau-I)^2=\tau I,
 \qquad \det S_\tau=1-\tau.
\]
The rank and trace are constant while the attachment deforms.
The algebraic families \eqref{cycen:eq:log-twist} and
\eqref{cycen:eq:formal-attaching} construct operations on
ordinary probe algebras. A Hall coproduct and factorization
sewing require their own maps.

\section{Half-braidings and central functors}
\label{cycen:sec:categorical}

The derived centre above is an algebra of operations. A
tensor category has a different centre, itself a braided
category. Its objects remember how one object exchanges
with every other object of the category.

Let $(\mathscr M,\otimes,\mathbf1)$ be a monoidal category.
An object of its \emph{Drinfeld centre} $\mathcal Z(\mathscr M)$
is a pair $(V,\beta)$, where $V\in\mathscr M$ and
\[
 \beta_W:V\otimes W\xrightarrow{\sim}W\otimes V
\]
is natural in $W$, equals the unit constraint at $\mathbf1$,
and obeys the half-braiding identity
\[
 \beta_{W\otimes W'}
   =(1_W\otimes\beta_{W'})(\beta_W\otimes1_{W'}).
\]
Associators are understood in this formula. Morphisms commute
with every $\beta_W$. Tensoring the half-braidings defines
the monoidal product, and $\beta_{V'}$ gives its braiding.

A central structure on a strong monoidal functor
$F:\mathscr B\to\mathscr M$ consists of half-braidings
on every $F(B)$, compatible with the tensor constraints.
If $\mathscr B$ is braided and the lift is to be braided,
its braiding must also map to the resulting half-braiding
between objects in the image. These structures are additional
to the monoidal constraints of $F$.

\begin{example}\label{cycen:ex:no-central-lift}
For a finite group $G$, let $\operatorname{Vec}_G$ be the
category of finite-dimensional $G$-graded vector spaces,
with $(V\otimes W)_g=\bigoplus_{ab=g}V_a\otimes W_b$.
Write $\delta_g$ for the one-dimensional object of degree $g$.
The symmetric category $\operatorname{Vec}_{C_2}$ maps
strongly monoidally to $\operatorname{Vec}_{S_3}$ by sending
the nontrivial degree to $s=(12)$. This functor has no
central lift. Indeed, a half-braiding of $\delta_s$ with
$\delta_h$, $h=(123)$, would identify objects of degrees
$sh$ and $hs$. These degrees differ, so every graded map
between them is zero.
\end{example}

Thus the centre construction does not provide a right adjoint
from arbitrary monoidal categories to braided categories with
all strong monoidal functors as morphisms. There is a genuine
adjunction for the forgetful functor from a fixed centre.
The finite-group case computes it completely.

\begin{proposition}\label{cycen:prop:group-centre}
The category $\mathcal Z(\operatorname{Vec}_G)$ consists of
$G$-graded vector spaces $V$ with a $G$-action $\rho$ such that
\[
 \rho(h)V_g\subset V_{hgh^{-1}}.
\]
The associated half-braiding on homogeneous elements is
\begin{equation}\label{cycen:eq:group-half-braiding}
 \beta(v\otimes w_h)=w_h\otimes\rho(h^{-1})v.
\end{equation}
The forgetful functor $U$ has right adjoint
\begin{equation}\label{cycen:eq:group-adjoint}
 I(X)=\bigoplus_{a\in G}(a,X),\qquad
 \deg(a,x_g)=aga^{-1},\quad
 \rho(h)(a,x)=(ha,x).
\end{equation}
\end{proposition}
\begin{proof}
The grading condition makes \eqref{cycen:eq:group-half-braiding}
degree-preserving. Exchanges first with degree $h$ and then
with degree $k$ apply
$\rho(k^{-1})\rho(h^{-1})=\rho((hk)^{-1})$; this is the
half-braiding identity. Conversely, the exchanges with all
$\delta_h$ determine these operators. The unit and tensor
identities give their group law, and degree preservation gives
the conjugation condition. Naturality extends the formula
from graded lines to all objects.

The object \eqref{cycen:eq:group-adjoint} has the required
grading and action. Given a graded map $f:UV\to X$, define
\[
 \widetilde f(v)=\sum_{a\in G}
       (a,f(\rho(a^{-1})v)).
\]
Each summand has the degree of $v$. Reindexing $a=hb$
proves equivariance, so $\widetilde f$ is a central morphism.
Projection to the identity-labelled summand recovers $f$.
Conversely, equivariance determines every summand of a central
map from that projection. The two constructions are inverse
and natural, proving
$\operatorname{Hom}_{\mathcal Z}(V,I(X))
 =\operatorname{Hom}_{\operatorname{Vec}_G}(UV,X)$.
\end{proof}

For $G=S_3$ and $X=\delta_{(12)}$, the object $I(X)$ has
dimension six. Each of the three conjugate transposition
degrees occurs twice. The adjoint adds these conjugate
components; it does not place a half-braiding on the original
one-dimensional object.

\begin{proposition}\label{cycen:prop:transport}
A strong monoidal equivalence
$F:\mathscr M\xrightarrow{\sim}\mathscr N$ induces a
braided equivalence
$\mathcal Z(\mathscr M)\xrightarrow{\sim}\mathcal Z(\mathscr N)$.
\end{proposition}
\begin{proof}
Choose a monoidal inverse $G$ and counit
$\varepsilon:FG\xrightarrow{\sim}\operatorname{Id}_{\mathscr N}$.
If $\phi$ denotes the tensor constraint of $F$, transport a
half-braiding by the composite
\[
 \begin{split}
 F(V)\otimes W&\xrightarrow{1\otimes\varepsilon_W^{-1}}
 F(V)\otimes F(GW)\xrightarrow{\phi}F(V\otimes GW)\\
 &\xrightarrow{F(\beta_{GW})}F(GW\otimes V)
 \xrightarrow{\phi^{-1}}F(GW)\otimes F(V)
 \xrightarrow{\varepsilon_W\otimes1}W\otimes F(V).
 \end{split}
\]
The monoidal coherence equations and the original
half-braiding identity prove the new one. The same construction
with $G$ supplies the inverse; the unit and counit give its
central equivalences. Applying the formula to two central
objects shows preservation of their braiding.
\end{proof}

A derived equivalence of geometric categories therefore
transports their ordinary derived centres by
Theorem~\ref{cycen:thm:morita}. To use
Proposition~\ref{cycen:prop:transport}, it must additionally
preserve the specified monoidal products. An equivalence of
underlying categories supplies no Hall tensor constraint.

\section{The internal centre for convolution}
\label{cycen:sec:convolution}

The coefficient category of the ordered collision schemes
carries a symmetric disjointness convolution. We now apply
the derived-centre construction internally, retaining the
coefficient object on every configuration.

Let $(\mathscr C_R,\star,K_R)$ be the presentable stable
coefficient category of Chapter~\ref{cycuts:ch}, extended
over a commutative $k$-algebra $R$. Its convolution preserves
colimits in both variables, and $K_R$ is supported at the
empty configuration. For an associative algebra $A$ in this
category, define $Z_\star^{\mathrm{int}}(A)$ by
\begin{equation}\label{cycen:eq:internal-centre}
 \operatorname{Map}_{\mathscr C_R}
       (V,Z_\star^{\mathrm{int}}(A))
 \simeq
 \operatorname{Map}_{A\text{-}A}(V\star A,A).
\end{equation}
The right side uses the two actions of $A$ on itself. Its
representability follows from presentability and the
colimit-preserving coefficient action.

An action on $A$ through convolution is a left-module
structure over an $E_2$ algebra $F$ in the category of
associative algebras of $\mathscr C_R$. Thus
$F\star A\to A$ is an algebra map with unit, multiplication
and action coherences. Morphisms of actions preserve all
these data.

\begin{theorem}\label{cycen:thm:convolution}
The internal centre is an $E_2$ algebra, and its evaluation
is final among convolution actions on $A$. Composition with
evaluation identifies maps of $E_2$ algebras
$F\to Z_\star^{\mathrm{int}}(A)$ with such actions of $F$.
Evaluation at the unit of $A$ gives an algebra map
$Z_\star^{\mathrm{int}}(A)\to A$.
\end{theorem}
\begin{proof}
The centre theorem for a presentable symmetric monoidal
category whose tensor preserves colimits constructs the final
action \cite{LurieHA}*{Corollary~5.3.1.15}. Its centralizer
of the identity is the morphism object of $A$ as a bimodule
\cite{LurieHA}*{Proposition~5.3.1.8 and Theorem~5.3.1.30}.
That morphism object is precisely
\eqref{cycen:eq:internal-centre}. The two compatible
compositions make it an associative algebra in associative
algebras, hence an $E_2$ algebra. Finally, compose evaluation
with $Z_\star^{\mathrm{int}}(A)\star K_R\to
Z_\star^{\mathrm{int}}(A)\star A$. Both arrows are algebra maps.
\end{proof}

For a left $A$-module $M$, tensoring evaluation over $A$
gives its characteristic action
\begin{equation}\label{cycen:eq:module-action}
 Z_\star^{\mathrm{int}}(A)\star M
 \simeq (Z_\star^{\mathrm{int}}(A)\star A)
                \star_A^{\mathbb L}M
 \longrightarrow A\star_A^{\mathbb L}M\simeq M.
\end{equation}
It is natural in $M$, including all homotopies. A
$\mathscr C_R$-linear equivalence of module categories
intertwines these actions and identifies their internal
centres as $E_2$ algebras. The proof of
Theorem~\ref{cycen:thm:morita} applies with $\star$ in place
of tensor over $k$: the coefficient-linear endofunctors are
bimodules and the identity is $A$. The coefficient linearity
is necessary to retain the internal object.

For $A=K_R\oplus a$ with positive augmentation ideal $a$,
let $[V,W]_\star$ denote the internal Hom for convolution.
The two-sided bar gives
\begin{equation}\label{cycen:eq:internal-bar}
 Z_\star^{\mathrm{int}}(A)
 \simeq
 \left(\prod_{q\geq0}[(sa)^{\star q},A]_\star,
                             \partial+[\tau,-]\right),
\end{equation}
where $\tau:BA\to A$ is the universal twisting cochain,
$\partial f=d_Af-(-1)^{|f|}fd_{BA}$, and
$f\smile g=\mu_A(f\star g)\Delta_{BA}$.
The adjacent bar products contribute to $\partial$;
the two outer module actions give the commutator with $\tau$.
The Maurer--Cartan equation for $\tau$ makes this differential
square zero. Cup and braces are those of
\eqref{cycen:eq:brace-sign}--\eqref{cycen:eq:operations},
using evaluation and the symmetry of convolution.

The product in \eqref{cycen:eq:internal-bar} can be infinite
even on one fixed output configuration: internal Hom from
a high convolution power can have a nonzero empty coefficient.
Define the scalar centre by
\[
 Z_\star^{\mathrm{sc}}(A)
  =R\operatorname{Hom}_{\mathscr C_R}
              (K_R,Z_\star^{\mathrm{int}}(A)).
\]
It retains only this scalar coefficient of the internal object.

For the square-zero algebra $K\oplus a$, the cochains with
values in $a$ have only their internal differential. The
positive-to-empty mapping objects vanish, so
\[
 Z_\star^{\mathrm{int}}(K\oplus a)
   \simeq K\oplus\prod_{q\geq0}[a^{\star q},a]_\star[-q].
\]
Their cup products vanish, but their braces need not. If
$\delta$ is the identity of $a$, regarded as a degree-one
Hochschild class, then any $a$-valued arity-$q$ cochain obeys
\begin{equation}\label{cycen:eq:arity-bracket}
                 [\delta,f]_{\mathrm G}=(1-q)f.
\end{equation}
The output insertion gives $f$ and the $q$ input insertions
give $qf$; $s\delta$ has degree zero, so no insertion sign occurs.

For $a=i_{\mathrm{ord}}\omega_{\mathbb A^1}$, every binary
reduced word gives the operation complex $k\oplus k[-1]$.
Indeed, applying derived Hom into the diagonal module to the
localization triangle gives a triangle with other terms
$k[-1]$ and $k[1]$. The map between them is zero over $k$;
the remaining complex is $k\oplus k[-1]$, with the residue
in degree zero. The middle-term calculation uses the one-variable
Spencer resolution: derivatives act freely on the delta module,
leaving one degree-one cokernel. The Hochschild shift by two
therefore gives scalar-centre classes $\mu_w,\nu_w$ of degrees
two and three, with
\[
 [\delta,\mu_w]_{\mathrm G}=-\mu_w,
 \qquad [\delta,\nu_w]_{\mathrm G}=-\nu_w.
\]
Thus even this square-zero example requires the brace operations.

\section{The ambient transverse endomorphism coalgebra}
\label{cycen:sec:transverse}

Convolution permits disjoint position sets. A transverse
bulk input may share every position of its boundary input.
The corresponding tensor is the diagonal tensor $\otimes^!$.
Let $\mathsf{Coalg}_+$ be the category of coaugmented counital
convolution coalgebras whose empty coefficient is $k$.
Its diagonal unit is the dualizing coalgebra $\mathbb J$,
with a vacuum at every configuration. Its reduced part is
positive. The empty convolution unit remains $K$.

A \emph{transverse monoid} is an associative algebra object
of $(\mathsf{Coalg}_+,\otimes^!,\mathbb J)$. Its product,
unit and action on a coalgebra $C$ have the forms
\[
 F\otimes^!F\longrightarrow F,
 \qquad \mathbb J\longrightarrow F,
 \qquad F\otimes^!C\longrightarrow C.
\]
They are coalgebra maps with coherent monoid or module laws.
In particular the product intertwines the coproduct through
the common-decomposition interchange of the two tensors.

\begin{theorem}\label{cycen:thm:ambient}
The category $\mathsf{Coalg}_+$ is closed for $\otimes^!$.
Write $[\![C,D]\!]$ for its internal mapping coalgebra. Then
\[
 \mathcal E(C)=[\![C,C]\!]
\]
is a transverse monoid, and evaluation classifies every
transverse monoid action on $C$:
\begin{equation}\label{cycen:eq:ambient-action}
 \operatorname{Map}_{\mathrm{Mon}}(F,\mathcal E(C))
 \simeq\{\text{coherent transverse actions of }F\text{ on }C\}.
\end{equation}
At one position,
\[
 \mathcal E(C)_{\{1\}}
 \simeq\underline{\operatorname{Hom}}_X^!
                    (C_{\{1\}},C_{\{1\}}).
\]
No factorization condition on $\mathcal E(C)$ is asserted.
\end{theorem}
\begin{proof}
Associative bar--cobar identifies $\mathsf{Coalg}_+$ with
augmented convolution algebras having positive augmentation
ideal. The latter form a presentable category. Limits of
nonunital algebras in the opposite coefficient category are
computed on underlying objects, so coalgebra colimits are
computed on reduced parts. The diagonal tensor preserves
these colimits componentwise. The adjoint functor theorem
therefore gives the mapping coalgebra.

An internal endomorphism object has composition and an identity.
The enriched endomorphism universal property identifies its
algebra maps with module structures, including their coherences
\cite{LurieHA}*{Corollary~4.7.1.40}. This gives
\eqref{cycen:eq:ambient-action}.

For the last statement, test the pointwise Hom on objects
supported on the small-diagonal branches. Local cohomology on
those branches is $i_{\mathrm{ord}}(C_{\{1\}})$, and the
projection formula reduces the adjunction to the one-position
mapping object. Bar words of length at least two vanish at
one position. The suspension in bar cancels the desuspension
in cobar, leaving the displayed endomorphism object.
\end{proof}

There is also an explicit cochain description. For reduced
parts $c$ of $C$ and $a$ of an augmented convolution algebra
$A$, put
\[
 \mathcal H(C,A)=\Omega[\![C,BA]\!].
\]
Its augmentation ideal is
$\underline{\operatorname{Hom}}^!(c,a)$, and on strict
derived representatives its product is determined by
\begin{equation}\label{cycen:eq:pointwise-cup}
 \operatorname{ev}(\cup\otimes^!1)
 =\mu_a(\operatorname{ev}\star\operatorname{ev})
              \zeta(1\otimes^!\overline\Delta_C).
\end{equation}
Here $\zeta$ projects onto common ordered decompositions.
The differential is the mapping differential, without a
twist by a distinguished cochain.

To verify the assertion, test on a positive coefficient $V$.
The bar of the free tensor algebra $T_\star V$ is equivalent
to the primitive coalgebra $K\oplus sV$. Its diagonal product
with $C$ is primitive on $sV\otimes^!c$, so free-algebra
adjunction gives
\[
 \operatorname{Map}(V,\overline{\mathcal H(C,A)})
 \simeq\operatorname{Map}(V\otimes^!c,a).
\]
The product \eqref{cycen:eq:pointwise-cup} is associative by
coassociativity, associativity and the coherent interchange.
Compose the bar twisting cochain of this cup algebra with
evaluation. The resulting twisting cochain defines a map
from its bar tensored with $C$ to $BA$; its adjoint induces
the identity on the displayed augmentation ideal. It is
therefore an equivalence of algebras. Taking $A=\Omega C$
gives
\[
 \Omega\mathcal E(C)\simeq\mathcal H(C,\Omega C).
\]
The identity of $C$ specifies its transverse unit.

The two action tensors cannot be interchanged in these
statements. For a positive coefficient $V$,
\[
 (K\star V)_{\{1\}}=V_{\{1\}},\qquad
 (K\otimes^!V)_{\{1\}}=0.
\]
The common-decomposition interchange has a section, but its
other composite is a projection, not an equivalence. The
convolution centre has the twisted differential in
\eqref{cycen:eq:internal-bar}; the transverse endomorphism
coalgebra has the untwisted cochains above. A comparison
must specify its tensor and its evaluation map.

\section{A Heisenberg monoid and its augmentation obstruction}
\label{cycen:sec:heisenberg}

The ambient construction does not require a chosen
augmentation to $\mathbb J$. A concrete factorizing monoid
acting on the Heisenberg boundary shows why such an
augmentation can fail even when its singleton cohomology
is an augmented algebra.

Fix $\kappa\in k^\times$, and let $V_\kappa$ be the
Heisenberg vertex algebra generated by $J$, with
$[J_\lambda J]=\kappa\lambda$. Adjoin an odd field $c$
of degree one, with regular brackets, and set
\[
 \mathcal Z_\kappa=V_\kappa\otimes\Lambda[c,Tc,T^2c,\ldots],
 \qquad dJ=\kappa Tc,\quad dc=0.
\]
The differential is the zero mode of ${:}cJ{:}$ and is a
vertex derivation. It squares to zero. All states are
algebraic polynomials in finitely many jets.

Let $E$ be the Weyl--Clifford vertex algebra on even fields
$t,p$ of degree zero and odd fields $\eta,\iota$ of degrees
$1,-1$, with
\begin{equation}\label{cycen:eq:quartet}
 [p_\lambda t]=1,\qquad [\iota_\lambda\eta]=1,\qquad
 Qt=\eta,\quad Q\iota=p,\quad Q\eta=Qp=0.
\end{equation}
The other generator brackets are zero or determined by
skewsymmetry. On $W=V_\kappa\otimes E$, the maps
\begin{equation}\label{cycen:eq:quartet-maps}
 i(J)=J+p,\qquad j(J)=J-\kappa Tt,\qquad j(c)=-\eta
\end{equation}
define vertex maps $i:V_\kappa\to W$ and
$j:\mathcal Z_\kappa\to W$ with commuting images.
Indeed,
\[
 [(J+p)_\lambda(J-\kappa Tt)]=\kappa\lambda-\kappa\lambda=0,
 \qquad Q(J-\kappa Tt)=-\kappa T\eta.
\]
Both image currents retain level $\kappa$.

The quartet is contractible away from its scalar vacuum.
On polynomial mode coordinates its differential pairs each
$t$ mode with its $\eta$ mode and each $\iota$ mode with
its $p$ mode. The reverse pairing, divided by the total
positive quartet degree, is a contraction. Thus $i$ is a
quasi-isomorphism. The commutant of $J+p$ has coordinates
\[
 v_n=x_n-\kappa n\tau_n,\quad \pi_n\quad(n\geq1),
 \qquad \bar\eta_m\quad(m\geq0),\quad\bar\iota_n\quad(n\geq1),
\]
Here $x_n=T^{n-1}J/(n-1)!$ and
$\pi_n=T^{n-1}p/(n-1)!$ for $n\geq1$,
whereas $\tau_m=T^mt/m!$ and
$\bar\eta_m=T^m\eta/m!$ for $m\geq0$.
Finally $\bar\iota_n=T^{n-1}\iota/(n-1)!$.
These are polynomial jet coordinates. Their differential is
$Qv_n=-\kappa n\bar\eta_n$ and $Q\bar\iota_n=\pi_n$.
Contracting the latter pairs shows that $j$ is a
quasi-isomorphism onto the commutant.

For $b$ quartets put $W_b=V_\kappa\otimes E^{\otimes b}$,
$K_b=J+\sum_{l=1}^bp_l$ and
$\mathcal Z^{(b)}=\operatorname{Com}(K_b,W_b)$.
These commutants have generators
$J-\kappa Tt_1$, $t_j-t_1$ for $j>1$, and all
$p_j,\eta_j,\iota_j$. For an order-preserving surjection
$\phi:[b']\twoheadrightarrow[b]$, define a merge
$m_\phi:W_b\to W_{b'}$ by
\[
 p_j\mapsto\sum_{\phi(l)=j}p_l,\quad
 \iota_j\mapsto\sum_{\phi(l)=j}\iota_l,\quad
 t_j\mapsto t_{\min\phi^{-1}(j)},\quad
 \eta_j\mapsto\eta_{\min\phi^{-1}(j)}.
\]
It fixes $J$ and carries $K_b$ to $K_{b'}$. The generator
brackets prove that it is a dg vertex map. The same polynomial
contractions show that it and its restriction to commutants
are quasi-isomorphisms. Taking the least element of a fibre
commutes with composition, so these merges form a diagram.

The shift $r_a:W_b\to W_{a+b}$ sends $J$ to $K_a$ and
the $j$th quartet to the $(a+j)$th. The products
\[
 \mathcal Z^{(a)}\otimes W_b\longrightarrow W_{a+b},
 \qquad z\otimes w\longmapsto{:}z\,r_a(w){:}
\]
are vertex maps because the two images commute. They restrict
to the commutants, are strictly associative and commute with
merges. Iteration sends the $i$th cochain factor to
\begin{equation}\label{cycen:eq:tower-currents}
 J^{(i)}\longmapsto J+\sum_{l<i}p_l-\kappa Tt_i,
 \qquad c^{(i)}\longmapsto-\eta_i,
\end{equation}
and the boundary current to $K_b$.

For a dg vertex algebra $B$ on $X=\mathbb A^1$, let
\[
 L_B=i_{\mathrm{ord}}(\omega_X\otimes B)
\]
be its main-diagonal Lie object for $\star$.
Here $i_{\mathrm{ord}}$ is the diagonal insertion of
Section~\ref{cycuts:sec:exchange}; the right differential-module
structure is induced by translation, and the bracket by the
vertex products. Write
\[
 C_B=\operatorname{CE}_\star(L_B)
     =K\oplus\bigoplus_{r\geq1}
                    \operatorname{Sym}_\star^r(L_B[1])
\]
for its Chevalley--Cousin factorization coalgebra.
The differential combines the internal differential with pairwise
bracket contractions, with the suspended Koszul signs; the
coproduct is the unshuffle coproduct. Each coefficient has
finitely many positive factors. Write $U_\star L_B$ for the
associative enveloping algebra of $L_B$.

\begin{proposition}\label{cycen:prop:heisenberg-monoid}
The coalgebra
$F_\kappa=C_{\mathcal Z_\kappa}$ has a coherent unital
transverse monoid structure and acts on $C_{V_\kappa}$.
Its products are the roofs
\[
 C_{\mathcal Z_\kappa}^{\otimes^! b}
 \longrightarrow C_{\mathcal Z^{(b)}}
 \xleftarrow{\sim}C_{\mathcal Z^{(1)}}
 \xleftarrow{\sim}C_{\mathcal Z_\kappa}.
\]
Its action is induced by the commuting maps in
\eqref{cycen:eq:quartet-maps} and the inverse of $C_i$.
\end{proposition}
\begin{proof}
For these free fields the factorization lattices are
$K$-flat: their PBW filtration has graded pieces equal to
symmetric powers of free principal-part modules, with a
linear differential. Tensoring lattices therefore computes
the derived diagonal tensor. The finite Chevalley--Cousin
comparison gives factorization and carries quasi-isomorphisms
to equivalences. It also carries polynomial homotopies to
homotopies. Hence $B\mapsto C_B$ preserves the tensor,
merges and strict products above.

The category of nonempty ordered finite sets and surjections
has a final singleton. Its nerve, and the nerves of its finite
products, are contractible. Localizing the diagram at its
merge equivalences therefore descends the products and action
to the singleton. The same argument for the two-coloured
algebra/module operad retains all their coherences: localization
of the associated cocartesian fibration localizes its fibres
\cite{cycen:HinichLocalization}*{Proposition~2.1.4}.

It remains to verify the unit. In two quartets, the merge
$mj$ and shift $r_1j$ are joined by the following two
polynomial homotopies, with parameter $u$ and forms on the left:
\[
 \begin{array}{c|cc}
  &J&c\\ \hline
  1&J-\kappa T((1-u)t_1+ut_2)
    &-(1-u)\eta_1-u\eta_2-du(t_2-t_1)\\
  2&J+up_1+du\,\iota_1-\kappa Tt_2&-\eta_2.
 \end{array}
\]
The first starts at $mj$; the second ends at $r_1j$.
They preserve brackets and differential and commute with
$K_2$. Right multiplication by the vacuum is the inclusion,
which agrees with $m$ on $j$; left multiplication is the
shift. Thus the vacuum is a quasi-unit. On the module its
action is an equivalence, since $r_1i_1=i_2$. Unitalization
of quasi-unital algebras and modules supplies the coherent
units \cite{LurieHA}*{Theorem~5.4.3.5 and
Proposition~5.4.3.14}. Transport along $C_j$ and $C_i$
gives the stated monoid and action.
\end{proof}

\subsection{The supported binary class}

The specified homotopies also determine a supported
operation. Put $Z=\mathcal Z_\kappa$ and
$Z'=\mathcal Z^{(2)}$. Its two embeddings are
\[
 \jmath_1(J)=U_1=J-\kappa Tt_1,\quad\jmath_1(c)=-\eta_1,
 \qquad
 \jmath_2(J)=U_1+p_1-\kappa T(t_2-t_1),\quad
 \jmath_2(c)=-\eta_2.
\]
Integrating the form components of the two homotopies gives
$h$ of degree $-1$, with
\[
 Qh+hd=\jmath_2-\jmath_1,\qquad
 h(J)=\iota_1,\quad h(c)=-(t_2-t_1).
\]
On a binary chart the unshifted coefficient is the fibre
of the chiral bracket
$\mu:j_*j^*(Z^r\boxtimes Z^r)\to\Delta_*Z^r$,
where $Z^r=\omega_X\otimes Z$. Its differential is
$d(\sigma,\zeta)=(d\sigma,\mu\sigma-d\zeta)$.
The two products of vacuum insertions are
$(\jmath_1\boxtimes\jmath_2,0)$ and
$(\jmath_2\boxtimes\jmath_1,0)$.

The external homotopy
\[
 H(a\boxtimes b)=-h(a)\boxtimes\jmath_2(b)
               +(-1)^{|a|}\jmath_2(a)\boxtimes h(b)
\]
has differential their difference in the first component.
Subtracting its boundary in the fibre leaves $(0,-\mu H)$.
Its Fourier expression is therefore
\begin{equation}\label{cycen:eq:roof-bracket}
 \beta(a,b)=[h(a)_\lambda\jmath_2(b)]
                 -(-1)^{|a|}[\jmath_2(a)_\lambda h(b)].
\end{equation}
This calculation specifies the overlap nullhomotopy as well
as the supported map.

Contract the extra quartet coordinates by the polynomial
contraction, obtaining $R:Z'\to Z$ with
$R\jmath_1=1$. The sign-adjusted pairing
$\{a_\lambda b\}=(-1)^{|a|}R\beta(a,b)$ has degree $-1$
and generator values
\[
 \{J_\lambda J\}=\{c_\lambda c\}=0,\qquad
 \{c_\lambda J\}=1,\quad\{J_\lambda c\}=-1.
\]
For example $[\jmath_2(J)_\lambda(t_2-t_1)]=-1$
and $[\iota_1{}_\lambda\eta_2]=0$, giving
$\beta(J,c)=-1$.

The underlying polynomial complex has coordinates
$x_n=J_{-n}\mathbf1$ for $n\geq1$ and
$\xi_n=T^nc/n!$ for $n\geq0$, with
\[
 dx_n=\kappa n\xi_n,\quad d\xi_n=0,\quad
 Tx_n=nx_{n+1},\quad T\xi_n=(n+1)\xi_{n+1}.
\]
Its cohomology is $\Lambda(\xi_0)$. The derivation
$K=\kappa^{-1}\sum_{n\geq0}x_{n+1}\partial_{\xi_n}$
satisfies $dK+Kd=T$. Thus the secondary value of the
specified pairing is
\begin{equation}\label{cycen:eq:secondary}
 -\{K\xi_0{}_\lambda\xi_0\}\big|_{\lambda=0}
                         =\kappa^{-1}.
\end{equation}
The homotopy $K$ refers to the polynomial complex; its
existence does not identify all its vertex operations with
those of a commutative vertex algebra.

To locate this value geometrically, record a moving
coefficient complex as a module over the de Rham algebra,
with differential $d+dz(\partial_z+T)$. The operator
$1-dzK$, with inverse $1+dzK$, changes this differential
to $d+dz\partial_z$. Indeed,
$[d+dz\partial_z,dzK]=-dzT$. Apply this change on the
two inputs and the diagonal output of the bracket.
The coefficient of $dz_1$ at Fourier parameter zero is
$(-1)^{|a|}\{Ka_\lambda b\}|_{\lambda=0}$.
For $a=b=\xi_0$ it is $\kappa^{-1}$.
The de Rham generator of the diagonal direct image is
$\varepsilon=[(dz_1-dz_2)\otimes1]$, the residue image
of $dz/z$, where $z=z_1-z_2$. Thus the supported lift has
class
\begin{equation}\label{cycen:eq:residue-lift}
 \kappa^{-1}(\xi_0\otimes\xi_0)^\vee\otimes\varepsilon.
\end{equation}

\subsection{Why the singleton augmentation does not extend}

Let $Y=\mathbb A^2$, $U=Y\setminus\Delta$ and
$\mathbb J_Y=\omega_Y[2]$. Put
$H_F=\Lambda(\xi_0)$ and
$P=(H_F\otimes H_F)\otimes\mathbb J_Y$.
The relative mapping group is
\begin{equation}\label{cycen:eq:relative-residue}
 \frac{H^{-1}R\operatorname{Hom}_U(P|_U,\mathbb J_U)}
 {\operatorname{im}H^{-1}R\operatorname{Hom}_Y(P,\mathbb J_Y)}
 \cong k(\xi_0\otimes\xi_0)^\vee[dz/z].
\end{equation}
This follows from $H^*_{\mathrm{dR}}(Y)=k$ and
$H^*_{\mathrm{dR}}(U)=k\oplus k[dz/z]$. The connecting
map is residue. If $N$ is the external nullhomotopy used
above and $L$ a chart nullhomotopy after applying a map
$f$ to $\mathbb J_Y$, the fibre convention gives
\[
 \Delta_*\epsilon_1\,\beta=-\partial[\,j^!fN-j^!L\,].
\]
Here $\epsilon_1:H_F\to k$ is the singleton augmentation.
Indeed, the iterated fibre with coordinates
$(\sigma,\zeta,N)$ maps to the supported complex by
$\zeta-\mu N$. Hence a pure overlap loop contributes
minus its residue. By \eqref{cycen:eq:residue-lift}, the
loop has residue $-\kappa^{-1}$ for every such $L$.

\begin{theorem}\label{cycen:thm:no-augmentation}
The transverse monoid $F_\kappa$ has no augmentation to
$\mathbb J$. Every transverse monoid receiving a monoid
map from $F_\kappa$ has the same obstruction, including
$\mathcal E(C_{V_\kappa})$. These assertions persist under
constant extension to any nonzero commutative $k$-algebra
concentrated in degree zero.
\end{theorem}
\begin{proof}
Suppose an augmentation exists. Use the roof model $Z'$,
and let $f$ be its two-position coefficient. Its coproduct
comparison supplies a homotopy $H$ from the localized $f$
to the external product of singleton augmentations.
Its multiplicativity comparison supplies a chart homotopy
$M$, and compatibility of that homotopy with coproduct
supplies a degree-$-2$ homotopy $\Gamma$ on $U$.
All three belong to the coherent augmentation data.

Denote by $\mu$ the two-factor product in the roof, by
$v',v''$ the two vacuum insertions, and by $s$ their symmetry.
The difference
\[
 L=M(v'\otimes^!v'')-M(v''\otimes^!v')s
\]
is a chart nullhomotopy of the augmented product difference:
the two products in the commutative unit cancel.
Expand the coproduct-compatibility equation for $M$ on
these two insertions. Its terms involving $H$ applied to
the product difference cancel against the same terms in
$fN$. The remaining loop is a chart restriction, a
$\Gamma$ boundary, and products of an insertion of $H$
with a closed chart map.

Each latter term is again a chart restriction plus a
boundary. On a single vacuum insertion, the possible
residue of $H$ belongs to
$\operatorname{Hom}^{-2}(H_F,k)[dz/z]=0$.
The degree-zero restriction on its boundary is injective,
so that boundary has a chart primitive. Subtracting it
leaves a closed insertion of $H$ whose only possible
non-chart component is the vanishing residue just described.
Multiplication by the other closed factor preserves exactness.
Thus the loop of this particular $L$ is zero in
\eqref{cycen:eq:relative-residue}. This contradicts the
value $-\kappa^{-1}$ computed above.

An augmentation of a receiving monoid would restrict to
one of $F_\kappa$, proving the inheritance assertion.
Its action gives the map to $\mathcal E(C_{V_\kappa})$
by Theorem~\ref{cycen:thm:ambient}. Over a degree-zero
coefficient algebra $R\ne0$, the same finite free mapping
calculation replaces $H_F$ by $R\oplus R\xi_0$ and the
residue by $-\kappa^{-1}1_R\ne0$. The vanishing unary
residue and the cancellation argument remain valid.
\end{proof}

A module whose underlying coalgebra is $\mathbb J$ is
precisely an augmentation: the endomorphism object of the
monoidal unit is the unit. Therefore neither left nor right
unit modules exist for these receiving monoids. The relative
bar with both endpoints equal to this fixed unit cannot be
formed. Other endpoint modules remain possible. Already
$M_2(k)$ has no one-dimensional unit module, but its row
and column modules have a well-defined relative tensor
product. Likewise a vacuum family varying through collisions
requires its own module and comparison data.

\section{Factorizing monoid actions}
\label{cycen:sec:native}

We now impose factorization on the acting coalgebra, while
retaining the transverse tensor and every action coherence.
The resulting universal problem can fail even for an ordinary
vertex algebra. The obstruction comes from a submonoid whose
action is trivial although its multiplication with another
submonoid leaves a nonzero binary operation.

For this section work over $\mathbb C$ on $X=\mathbb A^1$.
Let $Q$ be a unital vertex algebra in cohomological degree zero,
with algebraic state space and translation $T$. A
\emph{vertex derivation} is a linear map $D:Q\to Q$ such that
\[
 D\mathbf1=0,\qquad DT=TD,\qquad
 D(x_{(n)}y)=(Dx)_{(n)}y+x_{(n)}Dy
 \quad(n\in\mathbb Z).
\]
Write $C_Q$ for its Chevalley--Cousin factorization coalgebra
on the full ordered collision carrier, and put
$\mathfrak g=L_Q$, $\ell=L_{\mathbb C}$ and
$\mathbb J=C_{\mathbb C}$. The singleton coefficient of
$\mathfrak g$ is the right differential module
$\omega_X\otimes Q$, in degree zero.

The category of \emph{native monoid actions on $C_Q$} has
connected, coaugmented, counital factorizing coalgebras as
actors. Each actor has a coherent unital monoid structure for
$\otimes^!$ and a coherent action on $C_Q$. Morphisms preserve
the monoid and action structures, including their comparison
paths. Neither an augmentation to $\mathbb J$ nor a
propagation condition is imposed. Receiving actors may have
unbounded derived coefficients. A native monoid centre would
be a final object of this category.

\begin{theorem}\label{cycen:thm:native-obstruction}
If $Q$ admits a nonzero vertex derivation, its category of
native monoid actions has no final object. There are two
parallel morphisms of monoid actions which no receiving
factorizing monoid action can identify up to homotopy.
\end{theorem}

We construct the actors first. This also explains why a
calculation of formal coordinate maps alone would miss the
obstruction.

\subsection{Exterior actors and marked forms}

For a finite-dimensional vector space $W$ in degree zero, set
\[
 E_W=\Lambda(W[-1]),\qquad C_W=C_{E_W}.
\]
Thus the exterior generators have degree one; their
differential and translation are zero. The augmentation and
unit of $E_W$ give maps under and over $\mathbb J$.
A natural operation between these actors means a compatible
family formed from the ordinary chiral bracket, products in
the exterior coefficients, and the induced Chevalley and bar
maps. Only finite linear combinations are allowed in each
source weight and cochain degree. Completing over source
weights is harmless on a configuration of $m$ positions:
its exterior actor weight is at most $m\dim W$.

A natural map $C_W\to C_V$ has a marked matrix
\[
 A(t)\in\widehat{\operatorname{Sym}}(W^*)\otimes W^*\otimes V,
 \qquad \Phi_A(t)=A(t)t.
\]
The completion is in polynomial degree. One source root is
marked; exponential polarization supplies the factorials for
the remaining roots. The contraction $\Phi_A$ forgets that
mark. For example, $t\,ds-s\,dt$ is a nonzero marked form
with zero contraction.

\begin{lemma}\label{cycen:lem:marked}
Components of the natural mapping space under and over
$\mathbb J$ are the marked matrices $A(t)$, and each
component is contractible. Composition and exterior tensor
product are
\begin{equation}\label{cycen:eq:marked-law}
 A_{g f}(t)=A_g(\Phi_f(t))A_f(t),
 \qquad A_{f\otimes g}=A_f\oplus A_g.
\end{equation}
Different marked matrices remain different under geometric
homotopy, although this statement does not classify all
geometric maps.
\end{lemma}
\begin{proof}
Here is the natural mapping complex and its contraction.
Let $\mathfrak q_W$ be the Lie cobar of the reduced dual
coalgebra $(E_W^+)^*$, and adjoin the ordinary-input
generator $z$ of degree one, with $dz=\tfrac12[z,z]$.
Before taking its positive source-weight and
positive target-exterior summand, the
natural convolution algebra is the completed envelope of
\[
 \bigl(\operatorname{Lie}(z)\amalg\mathfrak q_W\bigr)
                         \otimes E_V.
\]
The coproduct in $E_W^*$ gives the cobar differential.
This description follows by resolving the input coalgebra:
connected rooted bracket expressions give the free Lie
generators, their ordered forests give its envelope, and
splitting an exterior input gives the cobar differential.
The relation $dz=z^2$ in the envelope accounts for the
ordinary bar input. Taking a positive target-exterior
factor imposes the relative condition under and over
$\mathbb J$.

The ordinary generator can be removed without changing
cohomology. In the augmentation ideal of
$(\mathbb C\langle z\rangle,dz=z^2)$, use
\[
 h(z^{2j})=z^{2j-1},\qquad h(z^{2j-1})=0
 \quad(j\geq1).
\]
The graded Leibniz rule gives $dh+hd=1$. In the free
product with $U(\mathfrak q_W)$, apply this contraction
to the first ordinary factor, with the sign of its prefix.
PBW symmetrization then proves the same assertion on
the Lie algebra, after tensoring with $E_V$, and after
taking its envelope.

The projection $\mathfrak q_W\to W^*$ kills brackets and
nonsingleton exterior labels and has abelian target.
To check it, write
$E_W=\Lambda(\alpha_1,\ldots,\alpha_r)$ and resolve the
augmentation module by
\[
 E_W\otimes\mathbb C[v_1,\ldots,v_r],
 \qquad |v_i|=0,\quad dv_i=\alpha_i.
\]
At positive combined weight $N$, the map
$N^{-1}\sum_i v_i\partial_{\alpha_i}$ contracts this
complex. Its dual derived Hom is concentrated in degree
zero and is $\operatorname{Sym}(W^*)$. The bar resolution
computes the same Hom; dualizing each finite weight
identifies it with $U(\mathfrak q_W)$. In degree zero,
its boundaries are exactly the commutators of the
singleton generators. PBW gives the asserted Lie
quasi-isomorphism.

Consequently the natural convolution algebra projects
surjectively, with weightwise acyclic kernel, to
\[
 \mathcal B(W,V)=
 \left[\widehat{\operatorname{Sym}}(W^*\otimes E_V)
                                      \right]_{E_V^+}.
\]
The subscript denotes the ideal generated by
$W^*\otimes E_V^+$. It has zero differential and no
terms in degrees at most zero. Its degree-one part is
precisely the space of marked matrices.

Lift the Maurer--Cartan equation one source weight at
a time. Its first curvature is a closed degree-two
kernel element, removed by a degree-one primitive.
The first difference of two lifts is removed by a
degree-zero primitive. Tensoring the same acyclic
kernel with polynomial forms relative to the boundary
of a simplex supplies all parameterized lifts.
Thus every fibre is contractible. In $\mathcal B$ a
degree-one Maurer--Cartan simplex is constant: its
square is zero by graded commutativity, and its
parameter differential vanishes.

To compute composition, keep independent ordinary root
variables and project their brackets to zero. The
marked source variable has one continuation through a
composite, giving matrix multiplication; the unmarked
variables substitute by $\Phi_f$. This gives
\eqref{cycen:eq:marked-law}. Disjoint coefficient
blocks give the tensor formula. The same top-root
coefficient is unchanged by a geometric homotopy:
after the ordinary variables become independent, its
possible cobar boundary has zero top coefficient.
It therefore distinguishes unequal matrices without
asserting that every geometric map is natural.
\end{proof}

The preceding contractions involve no infinite sum of
states or normal derivatives. At fixed source weight
and cochain degree the number of forest slots is
bounded; at a fixed configuration only finitely many
source weights occur. These are the finiteness
conditions used below.

\subsection{A monoid with a trivially acting direction}

Let $A=\mathbb C a\oplus\mathbb C b$ be the auxiliary
commutative algebra
\[
 a^2=a,\qquad ab=ba=b,\qquad b^2=0.
\]
Its unit $a$ is distinct from the vacuum of the
exterior vertex algebra. Put
\[
 F=C_A,\qquad G=C_{\mathbb C a},\qquad
 H=C_{\mathbb C b}.
\]
For $x\in A$, let $L_x$ be multiplication by $x$, and set
\[
 Q(x)=e^{L_x},\qquad
 M(x)=\sum_{n\geq0}\frac{L_x^n}{(n+1)!}.
\]
The inverse of $M$ exists as a formal series. Define
the multiplication of $F$ by the two marked blocks
\begin{equation}\label{cycen:eq:frame-product}
 \bigl(M(x+y)^{-1}M(x)\mid
              M(x+y)^{-1}Q(x)M(y)\bigr).
\end{equation}

\begin{lemma}\label{cycen:lem:frame-monoid}
This is a coherent unital transverse monoid under
and over $\mathbb J$. The inclusions
$i:G\to F$ and $j:H\to F$ are monoid maps.
The multiplication of $H$ is additive, and the
coordinate contraction of the multiplication of $F$
is $x+y$.
\end{lemma}
\begin{proof}
In the affine group with law
\[
 (x,U)(y,V)=(x+y,U+Q(x)V),
\]
associativity follows from $Q(x+y)=Q(x)Q(y)$.
The change of fibre coordinates $U=M(x)u$ gives
\eqref{cycen:eq:frame-product}, so its three
marked associativity identities and its unit
identities hold. In the algebra $A$,
\[
 M(x)x=e^x-a,\qquad
 e^{x+y}-a=(e^x-a)+e^x(e^y-a).
\]
Hence the diagonal $u=x$ is a homomorphism, which
proves the coordinate contraction. Both lines are
stable under these operators. On the $b$-line,
$L_{sb}b=0$, so both multiplication blocks are the
identity. Lemma~\ref{cycen:lem:marked} realizes the
identities and all their coherences. Its contractions
preserve the auxiliary grading
$\deg a=0$, $\deg b=1$, as does the multiplication.
\end{proof}

\begin{lemma}\label{cycen:lem:relative-action}
For a vertex derivation $D$ of $Q$, the monoid $F$
has a coherent action on $C_Q$ whose restriction
to $G$ is the trivial monoid action and whose
restriction to $H$ is the single-derivation action.
Both inclusions preserve these actions with their
specified comparisons.
\end{lemma}
\begin{proof}
First construct the operations with one formal
derivation. For $r$ copies of $F$, let $t_\nu,s_\nu$
be the coordinates dual to $a,b$. Give $D$ auxiliary
weight one, each $s_\nu$ weight minus one and each
$t_\nu$ weight zero. Prescribe the entire restriction
with no $b$-input to be trivial.

The scalar-bar and Koszul contractions in the preceding
proof apply also to differential polynomials in $D$.
They give the commutative relative model
\begin{equation}\label{cycen:eq:relative-derivation}
 \left[(s_1,\ldots,s_r)
 \mathbb C[[t_1,s_1,\ldots,t_r,s_r]]
                       \otimes\Lambda(\mathbb C D[-1])
 \right]_{\mathrm{weight}\,0}.
\end{equation}
The odd generator in this formula is the suspended
derivation, not the degree-zero endomorphism of $Q$.
The differential is zero. There are no nonpositive
degrees: a scalar monomial with a positive $s$-count
has negative auxiliary weight. The degree-one part is
exactly
\[
       \sum_{\nu=1}^r s_\nu f_\nu(t_1,\ldots,t_r)D.
\]
The projections to this model have the same
weightwise acyclic kernels as before.

Deleting an actor or deleting all its $a$-generators
commutes with the projection deleting the $b$-generators.
Their inclusions split the projections.
Inclusion--exclusion therefore extends compatible data
on the union of these faces; the remaining kernel is
a sum of multigraded acyclic kernels. The lifting
argument works relative to the unit faces, the
trivial $G$-restriction and a prescribed $H$-action.

Choose the class $sD$ for one copy. Pullback along the
monoid product is $(s_1+s_2)D$, since its coordinate
contraction is addition. Composition has the same
class: the formal operator identity
$e^{s_1D}e^{s_2D}=e^{(s_1+s_2)D}$ involves one
derivation. Relative contractibility supplies the
comparison path. Every higher module associahedron
has its boundary in the component with class
$(s_1+\cdots+s_r)D$. The same relative lifting fills
it, including all parameterized comparisons.

To evaluate these operations on the whole collision
carrier, the functor
$\mathcal T_Q(B)=B\otimes^!C_Q$ is strong symmetric
monoidal for $\star$. On each ordered disjoint
decomposition its comparison is the factorization
isomorphism of $C_Q$ followed by open projection.
Iteration gives the same ordered partition and
coproduct, and hence preserves collisions and normal
transfers. On the main diagonal it sends $\ell$ to
$\mathfrak g$ with its bracket. In particular
\[
 C_W\otimes^!C_Q
 \simeq\operatorname{CE}_\star(\mathfrak g\otimes E_W).
\]
These comparisons commute with envelopes and bars.
The vertex automorphism $1+uD$ over
$\mathbb C[u]/(u^2)$ carries $D$ to $D\otimes1$
through them, so the differential polynomials and
their homotopies evaluate to the required action.
Exterior degree bounds the total actor weight;
derivative order equals its $b$-count. Every
evaluated operation is finite on its specified input.
\end{proof}

The requirement that the whole $G$-action be trivial
is used here. The zero auxiliary degree of $a$ alone
would not remove scalar homotopies.

\subsection{A commutator with a nonzero marked form}

Let $S=C_{\mathbb C e_1\oplus\mathbb C e_2}$. Map its
exterior generators into four copies of $F$ by
\[
 e_1\longmapsto a^{(1)}-a^{(3)},\qquad
 e_2\longmapsto b^{(2)}-b^{(4)},
\]
and multiply in the displayed order. Denote the
resulting pointed coalgebra map by $w:S\to F$.
This map need not preserve a monoid structure on $S$.

\begin{lemma}\label{cycen:lem:commutator}
There is a natural pointed map $f:S\to H$ with
$w\simeq jf$. Its marked matrix is
\begin{equation}\label{cycen:eq:commutator-form}
 A_f(t,s)=\bigl(-s\chi(t)\mid t\chi(t)\bigr),
 \qquad \chi(t)=\sum_{n\geq0}\frac{t^n}{(n+1)!}.
\end{equation}
Its coordinate contraction and singleton components
are zero, but its first marked form is $t\,ds-s\,dt$.
\end{lemma}
\begin{proof}
Compute the fourfold product in the affine coordinates
of Lemma~\ref{cycen:lem:frame-monoid}, with base points
$ta,sb,-ta,-sb$. Their sum is zero, so the final
change of frame is the identity. If $dt,ds$ denote
the marked inputs, the fibre coordinate is
\[
 \begin{split}
 &M(ta)a\,dt+Q(ta)M(sb)b\,ds\\
 &\hspace{1em}-Q(ta+sb)M(-ta)a\,dt
                         -Q(sb)M(-sb)b\,ds.
 \end{split}
\]
The two $a$-components cancel. Using $b^2=0$,
the remaining component is
$b(-s\chi(t)\,dt+(e^t-1)\,ds)$.
This is \eqref{cycen:eq:commutator-form} in the
$b$-line. The natural classification supplies the
map and the homotopy. The other assertions follow
by contraction and by taking the lowest polynomial
degree.
\end{proof}

\subsection{Detection in an arbitrary receiving actor}

Fix $D\ne0$. Let $h:H\to Z$ be a receiving pointed
map of actions, and put $B=\Omega Z$. Write
$\iota:\ell\to B$ for its vacuum pointing and
$h_0$ for the pure singleton component of $h$.

\begin{lemma}\label{cycen:lem:singleton}
On every nonempty affine open $Y\subset X$,
\[
 [h_0]|_Y\ne0
 \quad\text{in}\quad
 H^1R\operatorname{Hom}_{Y}
              (\ell|_Y,B_{\{1\}}|_Y).
\]
No boundedness or holonomicity assumption on $B$ is needed.
\end{lemma}
\begin{proof}
We first examine a plain degree-zero coefficient class
$a:\ell|_Y\to B_{\{1\}}|_Y$. Its action gives an
endomorphism $M$ of $\mathfrak g|_Y$. We claim
\begin{equation}\label{cycen:eq:vacuum-module}
 M\mu_{\mathfrak g}
 =\mu_{\mathfrak g}(1\star M)
 =\mu_{\mathfrak g}(M\star1).
\end{equation}
Both ordered cobar products of $B$ are equivalences
on distinct roots. Indeed, the quotient by one product
summand is the cone of the other factorization
coproduct. Subtract the open exchange part of
$[\iota,a]$, with its nullhomotopy. Its image in
$U_\star\mathfrak g$ is zero there, because the
commutator in this envelope is the supported chiral
bracket.

The remaining map factors through
\[
 \operatorname{cofib}\bigl(j_!j^*(\ell\boxtimes\ell)
                    \longrightarrow
                         j_*j^*(\ell\boxtimes\ell)\bigr),
\]
where $j$ is the distinct-root inclusion. The source
is regular holonomic, so this $j_!$ is defined.
On either diagonal branch the cofiber has a
two-step filtration with successive terms $\ell[1]$
and $\ell$. Tensoring by $C_Q$ replaces them by
$\mathfrak g[1]$ and $\mathfrak g$. The PBW
envelope is in the right-module heart and its
extraordinary diagonal restriction is $\mathfrak g$.
Negative Ext kills the first term. The remaining
action has the form $N\mu_{\mathfrak g}$ and also
$\mu_{\mathfrak g}(1\star M)$. Inserting the vacuum
in the first slot gives $N=M$: the possible kernel
of the principal-part vacuum precomposition is
again a negative Ext between heart objects.
Skewsymmetry gives the other equality in
\eqref{cycen:eq:vacuum-module}.

Suppose now that $[h_0]|_Y=0$. A relative gauge in
the resolved pointed mapping complex removes
$h_0$, fixing $\mathbb J$. Its first actor-weight
term is then a degree-zero cocycle on one ordinary
input and one actor input. Choose $q\in Y$,
put the actor at $q$, and vary the ordinary input
on $Y'=Y\setminus\{q\}$.

The finite Spencer resolution of $\ell$ commutes
with tensoring by an arbitrary coefficient complex.
Algebraic K\"unneth therefore expresses this
operation as a finite sum of bidegrees $(r,-r)$.
After the action, the second factor lies in
$H^{-r}R\operatorname{Hom}_{\mathbb C}(Q,Q)$,
which vanishes unless $r=0$. The action path gives
\[
              D\otimes1_Q=\sum_\nu M_\nu\otimes N_\nu,
\]
where every $M_\nu$ satisfies
\eqref{cycen:eq:vacuum-module}. This equality holds
for the two-output classes: the possible
zero-column boundary vanishes by open adjunction,
$R\operatorname{Hom}(\mathfrak g,
             \operatorname{Sym}^2_\star\mathfrak g)=0$.

Choose a linear functional on the finite span of
$1,N_1,\ldots,N_m$ taking $1$ to $1$.
It expresses $D|_{Y'}$ as a linear combination of
the $M_\nu$. Thus $D$ is both a module map and a
derivation. Subtracting the derivation identity
from \eqref{cycen:eq:vacuum-module} and inserting
the vacuum in the second slot gives $D=0$.
Localization is faithful on the free algebraic
state module $\mathcal O_X\otimes Q$. This
contradicts the choice of $D$.
\end{proof}

\begin{lemma}\label{cycen:lem:binary-invariant}
The composite $hf:S\to Z$ is not homotopic under
$\mathbb J$ to the map through $\mathbb J$.
\end{lemma}
\begin{proof}
On the distinct-root slice
$(X\setminus\{q\})\times\{q\}$, its first binary class is
\begin{equation}\label{cycen:eq:binary-detector}
 [\iota]|_{X\setminus\{q\}}\otimes[h_0]_q
       -[h_0]|_{X\setminus\{q\}}\otimes[\iota]_q.
\end{equation}
The normalized fibre at $q$ is $i_q^!B[1]$.
The two terms have K\"unneth degrees $(0,1)$ and
$(1,0)$. The second is nonzero by the preceding
lemma and the unit identity: $[\iota]_q$ acts as
$1_Q$. Changing the cobar product order adds a
cobar boundary, so centrality of $\iota$ is not
being assumed.

We verify invariance under an arbitrary pointed
geometric homotopy. Represent it in the full
coefficient resolution by $\tau(u)+c(u)\,du$,
where $\tau$ and $c$ have degrees one and zero.
Both endpoint maps have zero pure singleton
components. On those inputs the interval equation
gives
\[
 \tau_{1,0}=\partial E,\qquad
 \tau_{0,1}=\partial F,\qquad
 E(u)=\int_0^u c_{1,0}(v)\,dv,\quad
 F(u)=\int_0^u c_{0,1}(v)\,dv.
\]
Here $E(0)=F(0)=0$ and
$\partial E(1)=\partial F(1)=0$.
Restriction to distinct roots kills the supported
source bracket. The binary equation becomes
\[
 \tau'_{1,1}=\partial c_{1,1}
                   +[\partial E,F']+[\partial F,E'].
\]
The brackets include the labelled external
products and the target multiplication. Since
$E,F$ have degree zero, the graded Leibniz rule gives
\[
 [\partial E,F']+[\partial F,E']
   =\partial[E,F']-\frac{d}{du}[E,\partial F].
\]
Integration makes the difference of the endpoint
binary components a boundary; the endpoint term
$[E,\partial F]|_0^1$ is zero. Their classes are
therefore equal under any such homotopy.
The map through $\mathbb J$ has zero binary class,
whereas \eqref{cycen:eq:binary-detector} is nonzero.
\end{proof}

\begin{proof}[Proof of Theorem~\ref{cycen:thm:native-obstruction}]
Let $i_0:G\to\mathbb J\to F$ be augmentation followed
by the unit. The relative action construction gives
two morphisms
\[
          (G,\mathrm{triv})\rightrightarrows(F,\alpha),
          \qquad i,\ i_0.
\]
If a receiving monoid-action map $k:F\to Z$ identified
$ki$ and $ki_0$, its multiplication coherences would
send the four-factor commutator $w$ through
$\mathbb J$: the first and third factors become
units, and the remaining $H$-factors cancel.
But $w\simeq jf$, and the map $h=kj$ is a receiving
pointed action map. Lemma~\ref{cycen:lem:binary-invariant}
forbids this homotopy. A final object would identify
the parallel maps after their map into it, so no
final object exists.
\end{proof}

For the Heisenberg vertex algebra, translation itself
is a nonzero vertex derivation, since $TJ\ne0$.
Thus the theorem applies to its ordinary factorization
coalgebra. The acting monoid of
Proposition~\ref{cycen:prop:heisenberg-monoid} and the
ambient universal monoid of
Theorem~\ref{cycen:thm:ambient} remain valid constructions.
The latter is not required to factorize. The parallel-pair
obstruction also differs from a claim that two objects
have no common recipient; no such claim is used or
deduced here.

A physical bulk reconstruction must consequently specify
its actor class, action tensor, local vacuum and morphisms
independently. Restricting that class may change the
universal problem, but then the restriction and the
comparison to physical observables are additional
mathematical data.

\section{Supported comparison equations}

The same open--closed decomposition decides whether a
map on the separated locus extends through a collision.
Let \(j:U\hookrightarrow X\) be an open embedding,
\(i:Z\hookrightarrow X\) its closed complement, and
work in a stable coefficient category with recollement
\[
 i_*i^!Q\longrightarrow Q\longrightarrow Rj_*j^*Q
                          \longrightarrow i_*i^!Q[1].
\]
For a source \(P\), a map \(f_U:j^*P\to j^*Q\)
has a connecting class
\[
 \operatorname{ob}(f_U)\in
              H^1R\operatorname{Hom}_X(P,i_*i^!Q).
\]

\begin{proposition}\label{cyp:prop:extension}
The map \(f_U\) extends to \(P\to Q\), with a specified
homotopy on restriction, exactly when
\(\operatorname{ob}(f_U)=0\). After choosing one extension,
the space of extensions is identified with
\(\operatorname{Map}_X(P,i_*i^!Q)\).
\end{proposition}

\begin{proof}
Apply \(R\operatorname{Hom}_X(P,-)\) to the recollement
triangle. Open adjunction identifies the third term with
\(R\operatorname{Hom}_U(j^*P,j^*Q)\). Exactness gives the
obstruction and the lift. Translation by a chosen lift
identifies the homotopy fiber over \(f_U\) with the fiber
over zero. This retains the restriction homotopy and all
higher homotopies of an extension.
\end{proof}

Extension is only the first equation for a comparison
preserving composition. To specify the higher equations,
let \(\mathsf A,\mathsf B\) be augmented homotopy
associative algebras over \(k\), with augmentation ideals
\(\overline{\mathsf A},\overline{\mathsf B}\).

Their bar coderivations \(b^{\mathsf A},b^{\mathsf B}\)
on the tensor coalgebras of the suspended ideals have
degree one and square zero. A coaugmentation-preserving
coalgebra map \(F\) of degree zero has Taylor components
\[
 f_n:(\overline{\mathsf A}[1])^{\otimes n}
                         \longrightarrow\overline{\mathsf B}[1],
                         \qquad n\geq1.
\]

All tensor substitutions use the Koszul signs on these
suspended complexes. Suppose the allowed derived Hom
complexes are closed under these substitutions, and let
\(\partial\) be their differential induced by
\(b_1^{\mathsf A}\) and \(b_1^{\mathsf B}\).
Projection of \(b^{\mathsf B}F=Fb^{\mathsf A}\) to one
output gives
\begin{equation}\label{cyp:eq:higher-comparison}
\begin{split}
 \partial f_n={}&
 \sum_{\substack{r+s+t=n\\s\geq2}}
 f_{r+1+t}(1^{\otimes r}\otimes b_s^{\mathsf A}
                                      \otimes1^{\otimes t})\\
 &-\sum_{\substack{\ell\geq2\\i_1+\cdots+i_\ell=n}}
 b_\ell^{\mathsf B}(f_{i_1}\otimes\cdots\otimes f_{i_\ell}).
\end{split}
\end{equation}
The first sum inserts one source operation. The second
runs over every ordered partition of the inputs into
blocks. In particular, at four inputs it contains
\(-b_2^{\mathsf B}(f_2\otimes f_2)\), even when the binary
defect has a primitive.

For an oriented Hall-to-field comparison, the endpoints,
orientations, supported mapping complexes and both unit
markings must first be specified. The regular generator
then reconstructs their maps. It neither chooses those
operations nor supplies the primitives required by
\eqref{cyp:eq:higher-comparison}.
